\documentclass{report}

\usepackage[dutch,english]{babel}   % load babel before spelatex to avoid
                                    % option clash!
\selectlanguage{english}
\usepackage[format=ogg,
            server=https://ctan.org/tex-archive/macros/latex/contrib/spelatex/Example
           ]
           {spelatex}

\newrobustcmd\CTAN{CTAN}
\spelmacad{CTAN}{see-tan}
\newrobustcmd\CPAN{CPAN}
\spelmacad{CPAN}{see-pan}

\title{\spelatex{} Example}
\author{Walter Daems and Paul Levrie}
\date{2025/01/31}
\setlength\parindent{0em}
\setlength\parskip{1ex}

\begin{document}

\maketitle

\chapter{Introduction}

\<<<{
  This file is just a simple showcase of the features of \spelatex{}.
  Below, you'll find examples of:
}

\begin{itemize}
  \spelitem{a simple equation}
  \spelitem{a more complex equation}
\end{itemize}

\chapter{Examples}
\section{A simple equation}
\label{eqn:simple}
\<<<{
  Consider the following simple definition of a polynomial function and
  check its spoken version by clicking on it.
}
\<<<[arealink]{
  \begin{equation}
    f(x) = x^{5}- x^4 + 7 x^3 + 3 x^2 - 8 x + 25
  \end{equation}
}
\<<<{
  This seems a simple equation, however, it is not so straightforward
  for an automated reader, to read it correctly.
}

\section{A more complex equation}
\newcommand\xx[2]{\ensuremath{#1_{#2}}}
\spelmacad{xx}[2]{#1 #2}

\label{eqn:complex}
\begin{spelchunk}
  For a lightray that hits the parabola at the point
  $P(t,9-\frac{t^2}{4})$, the reflected ray has slope $\tan 2\alpha$.
  Since the slope of the tangent to the parabola at $P$ is
  equal to $\tan\alpha = -\frac{t}{2}$, the equation of the
  reflected ray is given by
\end{spelchunk}
\begin{spelchunk}[arealink]
  \[
  y-9+\frac{t^2}{4} = -\frac{4t}{4-t^2} \cdot (x-t)
  \]
\end{spelchunk}

\section{Remark}
\<<<{
Instead of the \texttt{\textbackslash{}<<<} macro, one can also use the
spelchunk environment. We did this in the next sections.
}

\chapter{More advanced topics}

\selectlanguage{dutch}
\section{Een andere taal}
\begin{spelchunk}
\spelatex{} is ook volledig babel-actief, wat wil zeggen dat de
voorleesstem de geselecteerde taal zal volgen.
\end{spelchunk}

\begin{spelchunk}[arealink]
  \[
  y-9+\frac{t^2}{4} = -\frac{4t}{4-t^2} \cdot (x-t)
  \]
\end{spelchunk}

\selectlanguage{english}

\section{And some extras}
\subsection{Citations}
\begin{spelchunk}
Two excellent repositories are \CPAN{} \cite{CPAN} and \CTAN{} \cite{CTAN}.
\end{spelchunk}

\subsection{References to labels}
\begin{spelchunk}
Section~\ref{eqn:simple} contains an illustration of a simple
equation. For a more complex equation, we refer the user to
section~\ref{eqn:complex}.
\end{spelchunk}

\selectlanguage{dutch}

\bibliographystyle{alpha}

\begin{thebibliography}{99}

\bibitem{CTAN}
The Comprehensive \TeX{} Archive Network.
\newblock \url{http://www.ctan.org}.
\newblock online, accessed in August 2021.

\bibitem{CPAN}
The Comprehensive Perl Archive Network.
\newblock \url{http://www.cpan.org}.
\newblock online, accessed in August 2021.

\end{thebibliography}

\end{document}
